% unit: noether.p14.s05.title
\subsection*{5. Comparaison de la théorie arithmétique des fonctions algébriques avec la théorie géométrique. Concept d'invariance différent}

% unit: noether.p14.s05.par.001
Est \og{}invariant\fg{} tout concept caractérisé uniquement par le corps. Dans la théorie arithmétique, cette invariance est en général déjà donnée dans la définition, laquelle est énoncée par rapport à tous les éléments du corps, et non par rapport à une base quelconque ou à une variable distinguée; le concept de point indiqué plus haut en donne un exemple. En revanche, les autres théories partent d'une variable ou d'une base particulière et montrent après coup l'indépendance du concept à l'égard de ce choix particulier; l'invariance y devient une invariance vis-à-vis des transformations birationnelles.

% unit: noether.p14.s05.par.002
Supposons en effet que le corps soit obtenu, d'une part, par adjonction de l'élément algébrique \(s\) à \(K(z)\), et, d'autre part, par adjonction de \(\sigma\) à \(K(\xi)\), où \(f(z,s)=0\), respectivement \(F(\xi,\sigma)=0\), désignent les équations irréductibles de \(s\), respectivement de \(\sigma\); ou encore plus généralement par adjonction d'un nombre fini d'éléments à \(K(\eta)\), avec les équations irréductibles correspondantes
% unit: noether.p14.s05.eq.001
\[
g_1(\eta,\tau_1)=0,\qquad g_2(\eta,\tau_1,\tau_2)=0,\ldots .
\]
Alors, par définition, tous les éléments peuvent s'exprimer rationnellement au moyen de \(z\) et \(s\), aussi bien qu'au moyen de \(\xi\) et \(\sigma\), et aussi au moyen de \(\eta,\tau_1,\tau_2,\ldots\). En particulier, \(z\) et \(s\) sont donc exprimables rationnellement par \(\xi\) et \(\sigma\), et réciproquement; de même \(z\) et \(s\) le sont rationnellement par \(\eta,\tau_1,\tau_2,\ldots\), et réciproquement. La courbe plane \(f(z,s)=0\) passe donc, par \og{}transformation birationnelle\fg{}, dans la courbe plane \(F(\xi,\sigma)=0\), ou encore dans la courbe de l'espace des variables \(\eta,\tau_1,\tau_2,\ldots\) définie par \(g_1(\eta,\tau_1)=0; g_2(\eta,\tau_1,\tau_2)=0,\ldots\). L'invariance se manifeste donc ainsi: la définition énoncée pour une courbe spéciale demeure valable pour toutes les courbes qui en résultent par transformation birationnelle (dans un espace d'un nombre arbitraire de dimensions); pour toutes les courbes de la \og{}classe\fg{} au sens de Riemann. Ici, un \og{}point\fg{} est défini comme un système de valeurs associé \(z=a,s=b\), de sorte que \(f(a,b)=0\); par transformation birationnelle, ce point passe en général de nouveau dans un système de valeurs associé \(\alpha,\beta\) de \(F(\xi,\sigma)=0\), mais il peut, dans des cas particuliers, lorsque \(a,b\) était un point \og{}singulier\fg{}, correspondre à plusieurs \og{}points\fg{} de \(F(\xi,\sigma)=0\); et l'on peut toujours indiquer des courbes sur lesquelles un point \og{}singulier\fg{} de \(f(z,s)=0\) correspond à un nombre fini de points simples, séparés ou voisins (résolution des singularités).\footnote{Comparer par exemple Brill--Noether, 6e section. Des précisions sur les points \og{}singuliers\fg{} sont données dans le numéro suivant.} De même, un groupe de points passe de nouveau dans un groupe de points; et le nombre de ses points est invariant dès que l'on compte un point singulier pour autant de points que de points simples lui correspondent sous une transformation appropriée. Puisque toute courbe peut être transformée birationnellement en une courbe n'ayant que des points \og{}multiples ordinaires\fg{},\footnotemark[\value{footnote}] la théorie géométrique développe ses concepts seulement pour de telles courbes spéciales et montre ensuite l'invariance vis-à-vis de toute transformation birationnelle, même de celle qui conduit aux points singuliers les plus généraux.

% unit: noether.p14.s05.par.003
En particulier, comme courbe distinguée de la \og{}classe\fg{} --- le cas hyperelliptique excepté --- on peut introduire la \og{}courbe normale des \(\varphi\)\fg{} dans l'espace de \((p-1)\) dimensions, qui ne possède pas de points singuliers; ici les \(\varphi\) sont les \(p\) \og{}courbes adjointes d'ordre \((n-3)\)\fg{} pour \(f(z,s)=0\) rendu homogène. À une transformation birationnelle de \(f(z,s)=0\) en \(F(\xi,\sigma)=0\) correspond une transformation linéaire des \(\varphi\), de sorte qu'il s'agit ici d'invariance au sens projectif. La représentation de toutes les fonctions du corps comme fonctions rationnelles des \(\varphi\) est appelée \og{}représentation invariante\fg{} dans la théorie géométrique.\footnote{Comparer par exemple Brill--Noether, 8e section. Pour les fonctions algébriques de plusieurs variables, la question de savoir si l'invariance peut être ramenée à une invariance vis-à-vis des transformations linéaires n'a pas été traitée.}

% unit: noether.p14.s05.par.004
Le faisceau linéaire des \(\varphi\), ou encore des \og{}groupes spéciaux\fg{} qu'elles découpent, passe donc dans lui-même sous toute transformation birationnelle; il est par conséquent caractérisé uniquement par le corps. Il correspond à la classe différentielle dans la théorie arithmétique, où l'invariance se montre de nouveau directement dans la définition, laquelle est donnée par des propriétés relatives à chaque fonction du corps comme variable indépendante.

% unit: noether.p14.s06.title
\subsection*{6. Suite de la comparaison. Points singuliers}

% unit: noether.p14.s06.par.001
C'est à la différence de conception du concept de point que se rattache le fait que les \og{}points singuliers\fg{} des courbes, qui causent des difficultés particulières dans la théorie algébro-géométrique, n'apparaissent pas du tout dans la fondation de la théorie arithmétique; et, par suite, que les \og{}fonctions adjointes\fg{} ne sont pas requises pour construire la théorie arithmétique, comme c'est le cas dans la théorie géométrique, mais appartiennent à des parties spéciales et plus élevées de la théorie. D'autre part, les points de ramification n'apparaissent pas dans la fondation de la théorie géométrique; en fait, ils disparaissent aussi dans la représentation des intégrales sous la forme normale d'Aronhold (Hensel no 27, formule 80).

% unit: noether.p14.s06.par.002
La définition du point singulier donnée en 5 peut être formulée ainsi: \(f(z,s)=0\) (respectivement une courbe dans un espace de dimension supérieure) possède en \(z=a,s=b\) (respectivement \(z=a; s_i=b_i\)) un point singulier si ce système de valeurs est pris par plus d'un point de la surface de Riemann absolue (point au sens arithmétique), ou s'il est pris en un ou plusieurs points avec un ordre plus élevé. Dans ce dernier cas, on a affaire à un point du polygone de ramification de \(z\) aussi bien que de \(s\); cela correspond aux \og{}points consécutifs\fg{} de la théorie géométrique; le point devient un point de rebroussement (supérieur) de \(f(z,s)=0\) (respectivement de la courbe dans l'espace supérieur).

% unit: noether.p14.s06.par.003
Comme, par hypothèse, toute fonction arbitraire \(\eta\) du corps peut être représentée rationnellement au moyen de \(z\) et \(s\), tout point de la surface de Riemann absolue où \(z=a,s=b\) est obtenu, en vertu de l'isomorphie, en assignant à chaque fonction \(\eta\), représentée par \(z\) et \(s\), sa valeur pour \(z=a,s=b\). Pour que le point soit singulier, il doit donc exister au moins une fonction \(\eta\) qui, pour \(z=a,s=b\), prend plusieurs valeurs ou prend plusieurs fois un système de valeurs. Et cela, à cause de l'exprimabilité rationnelle au moyen de \(z\) et \(s\), ne peut avoir lieu que si, dans cette représentation, le numérateur et le dénominateur s'annulent simultanément pour \(z=a,s=b\). Il s'ensuit que la définition ci-dessus est identique à la définition usuelle, selon laquelle
% unit: noether.p14.s06.eq.001
\[
f,\quad \frac{\p f}{\p z},\quad \frac{\p f}{\p s}
\]
s'annulent pour \(z=a,s=b\). Car dans ce cas
% unit: noether.p14.s06.eq.002
\[
\eta=\frac{\frac{\p f}{\p z}}{\frac{\p f}{\p s}}
\]
est une fonction du genre indiqué qui est multivaluée pour \(z=a,s=b\), tandis que, dans le cas contraire, toute fonction, même si elle devient \(0/0\), ne prend qu'une seule valeur.

% unit: noether.p14.s06.par.004
Avec le concept arithmétique de point, le point singulier ne peut pas apparaître, puisque deux points sont considérés comme distincts dès qu'une seule fonction se voit assigner des systèmes de valeurs différents. Mais même dans la théorie des idéaux qui sert à fonder le concept arithmétique de point, le concept de point singulier n'intervient pas, précisément parce que l'on considère toujours simultanément la totalité de toutes les grandeurs entières du corps. Tous les points singuliers situés dans le fini de \(F(z,\delta)=0\), où \(\delta\) est algébriquement entier relativement à \(z\), sont en effet engendrés (d'après la deuxième définition et la fin du no 3) par l'\og{}idéal du point double\fg{} \(\frf\). Or, puisque, d'après le théorème fondamental, l'idéal de ramification \(\frD\) est le plus grand commun diviseur de tous les idéaux principaux \(F'(\delta)\), on peut, pour chaque idéal premier \(\frp\), indiquer un \(\delta\) tel que le \(\frf\) correspondant ne soit pas divisible par \(\frp\); par suite --- puisque \(\frf\) est le plus grand commun diviseur de \(F'(\delta)\) et de \(F'(z)\) --- au moins l'un des deux idéaux principaux \(F'(\delta)\) et \(F'(z)\) n'est pas divisible par \(\frp\). Il n'existe donc \emph{aucun système de valeurs pour lequel tous les \(F'(\delta)\) et \(F'(z)\) s'annulent simultanément}, et par conséquent le système de valeurs \(z=a,\delta_i=b_i\) n'est pris qu'en un seul point de la surface de Riemann absolue; la \emph{totalité des fonctions algébriques entières du corps} (que l'on peut interpréter comme une courbe dans l'espace d'une infinité de dimensions) \emph{ne possède aucun point singulier dans le fini}; or seuls les valeurs finies entrent en considération dans la théorie des idéaux.

% unit: noether.p14.s06.par.005
Mais les fonctions de base \(\omega_1,\ldots,\omega_n\) de toutes les grandeurs entières ne possèdent pas non plus de points singuliers dans le fini. D'après ce qui précède, tout point au sens arithmétique est déjà déterminé de façon univoque par un système de constantes assigné isomorphiquement à toutes les fonctions entières. Si donc une \og{}courbe spatiale\fg{} définie par certaines fonctions entières \(s_i\) possède, pour \(s_i=\alpha_i\), un point singulier dans le fini, il doit exister au moins une fonction entière \(\eta\) qui, pour \(s_i=\alpha_i\), prend plusieurs systèmes de valeurs. Si maintenant \(\omega_1,
\ldots,\omega_n\) désigne une base de toutes les grandeurs entières, alors, puisque \(\delta=a_1\omega_1+\cdots+a_n\omega_n\) avec des \(a_i\) entiers rationnels, à chaque système de valeurs des \(\omega\) ne correspond qu'un seul système de valeurs de \(\delta\); par conséquent la courbe spatiale définie par \(\omega_1,
\ldots,\omega_n\) ne peut posséder aucun point singulier dans le fini; et \emph{chaque point de la surface de Riemann absolue où \(z\) est fini est déjà défini de façon univoque par un système de constantes assigné isomorphiquement aux \(\omega\). Les fonctions de base sont précisément les fonctions \(\eta\) qui deviennent multivaluées pour \(s_i=\alpha_i\).} Par l'introduction de la base des grandeurs entières, par exemple à la place d'une base \(1,\delta,\ldots,\delta^{n-1}\), la difficulté des points singuliers est donc contournée.

% unit: noether.p14.s06.par.006
Dans la théorie géométrique, la tâche d'engendrer univoquement tous les points de la surface de Riemann absolue --- plus généralement, de former toutes les fonctions qui deviennent infinies en des points donnés --- est accomplie au moyen des formes adjointes et de leurs quotients. Une forme \(\psi\) est dite adjointe si \(\psi=0\) possède au moins un point \((i-1)\)-uple en chaque point \(i\)-uple de la courbe fondamentale \(f(z,s)=0\); un point singulier supérieur doit d'abord être résolu en points multiples ordinaires. Il a été montré plus haut que, pour la fonction la plus générale \(\eta=\frac{\psi}{\chi}\), le numérateur et le dénominateur doivent de toute façon s'annuler en un point singulier de \(f(z,s)=0\); qu'ils doivent être adjoints, cela ressort de l'examen des différentielles partout finies. Inversement, que les conditions d'adjonction soient également suffisantes pour former la fonction la plus générale, c'est ce que montre le \og{}théorème des résidus\fg{}, sur lequel on reviendra plus précisément. Les fonctions adjointes prennent donc la place des fonctions de base de la théorie arithmétique.

% unit: noether.p14.s06.par.007
Arithmétiquement, à une forme adjointe correspond le numérateur d'une fonction du corps divisible par l'\og{}idéal du point double\fg{} \(\frf\), dès que l'on suppose que les points singuliers --- ce qui peut toujours être obtenu par transformation --- sont tous situés dans le fini. De là se montre formellement la représentation des grandeurs de base \(\omega_1,
\ldots,\omega_n\) comme quotients de formes adjointes, et par suite le rapport entre les différentes façons de les représenter, comme suit:

% unit: noether.p14.s06.par.008
Soit \(R(z)\) le déterminant de substitution qui transforme une base \(1,\delta,\ldots,\delta^{n-1}\) en une base \(\omega_1,
\ldots,\omega_n\). Réciproquement, chaque \(\omega\) s'exprime alors comme quotient d'une fonction entière de \(z\) et \(\delta\) par \(R(z)\):
% unit: noether.p14.s06.eq.003
\[
\omega_i=\frac{G_i(z,\delta)}{R(z)}.
\]
Comme \(R(z)\omega_i\) appartient à l'anneau dérivé de \(\delta\), il résulte de la définition du conducteur que \(R(z)\) est divisible par \(\frf\) [de \((R(z))^2=\operatorname{Norm}\frf\), on ne peut conclure cette divisibilité que pour \((R(z))^2\)]. Puisque maintenant \(\omega_i\), en tant que fonction entière, reste finie pour tout \(z\) fini, la fonction numérateur \(G_i(z,\delta)\) doit aussi être divisible par \(\frf\); les numérateurs et les dénominateurs de tous les \(\omega_i\) sont donc adjoints. Cette représentation des \(\omega\) donne alors celle de toutes les fonctions du corps par quotients d'adjointes.

% unit: noether.p14.s06.par.009
D'après ce qui précède, la théorie géométrique peut être caractérisée comme la \emph{théorie de l'anneau dérivé d'une grandeur \(\delta\)} (ordre), par opposition à la théorie arithmétique, qui considère \emph{toutes les grandeurs entières} simultanément. Il est donc clair que le conducteur de l'anneau --- les points singuliers --- joue un rôle décisif. On peut encore indiquer d'autres analogies entre la théorie arithmétique des anneaux (ordres) et la théorie géométrique. Ainsi, la théorie géométrique ne compte nulle part les points d'intersection qui tombent dans les points singuliers lors de l'intersection avec des courbes adjointes, mais ne considère que les groupes de points qui en sont distincts. À cela correspond le fait que Dedekind\footnote{Über die Anzahl der Idealklassen in den verschiedenen Ordnungen eines endlichen Körpers. (Festschrift Braunschweig 1877); \S{} 3--5. (Les théorèmes énoncés dans ces paragraphes pour les nombres algébriques valent littéralement pour les fonctions algébriques d'une variable.)} ne considère que les idéaux d'anneau qui sont premiers au conducteur \(\frf\), et qu'il établit, pour ce domaine d'idéaux, toutes les lois de divisibilité, y compris la décomposition unique en idéaux premiers.

% unit: noether.p14.s06.par.010
Le \og{}théorème fondamental des fonctions algébriques\fg{}\footnote{M. Noether, Über einen Satz aus der Theorie der algebraischen Funktionen. \emph{Math. Ann.} vol. 6 (1873).} qui est à la base de la théorie géométrique, ou du moins une conséquence immédiate de celui-ci, trouve aussi, en un certain sens, un analogue dans la théorie arithmétique des anneaux. Ce théorème donne les conditions d'une représentabilité
% unit: noether.p14.s06.eq.004
\[
\psi(z,s)=A(z,s)f(z,s)+B(z,s)\varphi(z,s),
\]
où toutes les grandeurs qui interviennent sont entières rationnelles en \(s,z\) (plus généralement entières et homogènes en \(x_1,x_2,x_3\)). De ces conditions, il s'ensuit,\footnote{Brill--Noether, no 55.} en vertu de \(f(z,s)=0\), que le quotient \(\frac{\psi}{\varphi}\) est toujours égal à une fonction entière rationnelle \(B(z,s)\) lorsqu'il est adjoint à \(f\). À cela correspond le théorème du no 3, qui est à la base de toutes ces comparaisons: le conducteur \(\frf\) d'un anneau dérivé de \(\delta\) est égal à l'idéal du point double de \(f(z,\delta)=0\). Car, par la définition du conducteur, toute grandeur algébrique entière divisible par \(\frf\) appartient à l'anneau, et elle est donc entière et représentable rationnellement par \(z\) et \(\delta\); tandis que le théorème cité dit qu'une telle grandeur est une fonction adjointe.\footnote{On suppose ici constamment, à la suite de Dedekind, que \(\delta\) est une grandeur algébrique entière (ce que l'on peut toujours obtenir par transformation). Le cas correspondant plus exactement à la géométrie, \(f(z,s)=0\), où \(s\) peut aussi dépendre algébriquement de façon fractionnaire de \(z\), est traité par Hensel--Landsberg; comparer Hensel nos 26 et 29.}
